% not printed: handwritten call number "QA153 W35 v.1"; library stamps "MATH.-STAT. LIBRARY" and "M777593"
\begin{center}
{\Large \gesperrt{VORWORT.}}\\[1ex]
\rule{4em}{0.4pt}
\end{center}

Bei der Entwickelung, welche die Algebra in den letzten Jahrzehnten genommen hat, dürfte eine zusammenfassende Darstellung und Verknüpfung der verschiedenen theoretischen Betrachtungen und mannigfachen Anwendungen auch nach dem für seine Zeit trefflichen Lehrbuch von \gesperrt{Serret} nützlich sein.

Seit Jahren hege ich den Plan eines solchen Unternehmens, das ja gross und weitaussehend erschien, und mancherlei Vorarbeiten erforderte. Erst nachdem ich in Universitätsvorlesungen mehrmals das Gebiet im Ganzen durchwandert und einzelne Theile specieller behandelt hatte, entschloss ich mich, an die Ausführung des Werkes zu gehen, von dem jetzt der erste Band vollendet vorliegt.

Es war meine Absicht, ein Lehrbuch zu geben, das, ohne viel Vorkenntnisse vorauszusetzen, den Leser in die moderne Algebra einführen und auch zu den höheren und schwierigeren Partien hinführen sollte, in denen das Interesse an dem Gegenstande erst recht lebendig wird. Dabei sollten die erforderlichen Hülfsmittel, die elementaren sowohl als die höheren, aus dem Gange der Entwickelung selbst abgeleitet werden, um die Darstellung von anderen Lehrbüchern möglichst unabhängig zu machen.

Zwei Dinge sind es, die für die neueste Entwickelung der Algebra ganz besonders von Bedeutung geworden sind; das ist auf der einen Seite die immer mehr zur Herrschaft gelangende Gruppentheorie, deren ordnender und klärender Einfluss überall zu spüren ist, und sodann das Eingreifen der Zahlentheorie.
